\documentclass[10pt]{article}

\usepackage[letterpaper,margin=0.78in,headheight=14pt]{geometry}
\usepackage[T1]{fontenc}
\usepackage{lmodern}
\usepackage{microtype}
\usepackage{amsmath,amssymb,mathtools,amsthm}
\usepackage{booktabs,array,tabularx,multirow,graphicx}
\usepackage{xcolor}
\usepackage{hyperref}
\usepackage{enumitem}
\usepackage{fancyhdr}
\usepackage{listings}
\usepackage{titlesec}
\usepackage{lastpage}
\usepackage{needspace}

\widowpenalty=10000
\clubpenalty=10000
\displaywidowpenalty=10000

\definecolor{Navy}{HTML}{102D44}
\definecolor{Teal}{HTML}{0D716A}
\definecolor{Pale}{HTML}{EDF5F4}
\definecolor{Gray}{HTML}{52606D}
\definecolor{Red}{HTML}{9B2C2C}

\hypersetup{
  colorlinks=true,
  linkcolor=Navy,
  urlcolor=Teal,
  citecolor=Teal,
  pdftitle={Pair-Sum State Collapse in Rainbow and the Design of Rainbow 2},
  pdfauthor={Cris},
  pdfsubject={Chosen-Prefix Collisions, Finalization Bottlenecks, and the Hardening Ladder from Rainbow v3.7.1 to Rainbow 2},
  pdfkeywords={Rainbow; Rainbow 2; hash functions; collision attacks; chosen-prefix collisions; non-cryptographic hashing; state collapse; SMHasher3; AI-assisted cryptanalysis}
}
\urlstyle{same}

\pagestyle{fancy}
\fancyhf{}
\lhead{Cryptanalysis of Rainbow v3.7.1 \& Design of Rainbow 2}
\rhead{Cris}
\cfoot{\thepage\ / \pageref{LastPage}}
\renewcommand{\headrulewidth}{0.4pt}

\titleformat{\section}{\large\bfseries\color{Navy}}{\thesection}{0.55em}{}
\titleformat{\subsection}{\normalsize\bfseries\color{Teal}}{\thesubsection}{0.55em}{}
\titlespacing*{\section}{0pt}{1.25em}{0.35em}
\titlespacing*{\subsection}{0pt}{0.9em}{0.25em}

\setlist[itemize]{leftmargin=1.3em,itemsep=0.12em,topsep=0.2em}
\setlist[enumerate]{leftmargin=1.5em,itemsep=0.15em,topsep=0.2em}

\lstdefinestyle{code}{
  basicstyle=\ttfamily\footnotesize,
  backgroundcolor=\color{Pale},
  frame=single,
  rulecolor=\color{Teal!45},
  breaklines=true,
  columns=fullflexible,
  xleftmargin=0.25em,
  xrightmargin=0.25em,
  aboveskip=0.45em,
  belowskip=0.45em
}

\newtheorem{theorem}{Theorem}
\newtheorem{lemma}{Lemma}
\newtheorem{proposition}{Proposition}
\newtheorem{corollary}{Corollary}
\theoremstyle{definition}
\newtheorem{definition}{Definition}
\newtheorem{remark}{Remark}
\newtheorem{result}{Experimental Result}

\newcommand{\W}{\mathbb{Z}/2^{64}\mathbb{Z}}
\newcommand{\ROTR}{\operatorname{ROTR}}
\newcommand{\mixA}{\operatorname{mixA}}
\newcommand{\mixB}{\operatorname{mixB}}
\newcommand{\code}[1]{\texttt{#1}}

\begin{document}

\begin{center}
  {\LARGE\bfseries\color{Navy} Pair-Sum State Collapse in Rainbow and the Design of Rainbow~2}\par
  \vspace{0.3em}
  {\large Full-State Collisions, Finalization Bottlenecks, and a Speed-Preserving Hardening Ladder}\par
  \vspace{0.6em}
  {\large Cris}\par
  DOSAYGO Corporation\par
  \href{mailto:research@dosaygo.com}{\texttt{research@dosaygo.com}}
\end{center}

\noindent\textbf{AI-assisted research disclosure.}
The analysis reported here was conducted with substantial assistance from Anthropic's Claude (Opus~5) and Google DeepMind's Antigravity for hypothesis refinement, derivation, implementation, and experiment design. The human author selected the target, reviewed the evidence, and takes responsibility for the claims and the manuscript. All computational evidence is published as source, witnesses, and machine-readable artifacts; every collision and benchmark is replayed against the published implementations.

\begin{abstract}
Rainbow (v3.7.1) is a fast general-purpose hash with 64-, 128-, and 256-bit digests, documented by its author as non-cryptographic. In the first part of this paper, we show that its 256-bit internal state carries only 128 bits of collision-relevant information, and that those 128 bits decompose into two independently attackable 64-bit halves. Because each 16-byte block is injected as $h_0 \mathrel{-}= x$, $h_1 \mathrel{+}= x$, $h_2 \mathrel{+}= y$, $h_3 \mathrel{-}= y$, injection preserves the two \emph{pair sums} $\sigma_1 = h_0 + h_1$ and $\sigma_2 = h_2 + h_3$ while moving the state freely inside their coset. We prove that two states at the same block index can be driven to the \emph{identical} 256-bit state by one further message block, with zero search, if and only if their pair sums agree (Theorem~\ref{thm:pairsum}). Because $\mixA$ acts independently on $(h_0,h_1)$ and $(h_2,h_3)$, equalizing the two sums costs two 64-bit birthday searches rather than one $2^{128}$ search, yielding a same-seed collision in $2^{32.28}$ evaluations ($6.6$\,s on 8 cores) and a chosen-prefix collision in $2^{34.80}$ evaluations ($50$\,s) that collide simultaneously at 64, 128, and 256 bits. We further characterize three additional structural weaknesses (`Q18`--`Q22`): a 192-bit finalization bottleneck in which word $h_0''$ is $100\%$ dead in 64-bit mode and the $w=4$ toy finalizer reaches only $6.1\%$ of the state space; 2-adic differential filtration rays across odd-prime multiplications ($M_r = 11.01\times$); and a linear top-bit seed schedule in which flipping $\Delta\text{seed}=2^{63}$ alters only 4 bits of the 256-bit IV.

In the second part of this paper, we construct and evaluate a seven-rung hardening ladder (\code{R0}--\code{R6}) to eliminate all five classes of structural weaknesses (\code{W1}--\code{W5}) while preserving Rainbow's 4.6--4.8\,GiB/s bulk throughput. Combining algebraic verification with 32-core cloud \code{SMHasher3} testing on AMD EPYC Milan uncovered a subtle bit-0 parity cancellation ($(Y \oplus X) - X \equiv Y \pmod 2$) when post-multiply Feistel cross-XORs were paired with an unrotated XOR-subtract squeeze, which we prove and eliminate via the rotated four-lane squeeze $g = \ROTR(h_0 - h_1, 29) - h_2 - h_3$. Based on the complete algebraic and empirical dataset, we select rung \code{R5} as \textbf{Rainbow~2} (\code{rainbow2}): a separate hash family that eliminates every structural weakness identified in Rainbow~1, passes the full \code{SMHasher3} statistical suite at 64, 128, and 256 bits, and sustains $1.41$--$1.44$\,bytes/cycle ($4.60$--$4.69$\,GiB/s) bulk throughput.
\end{abstract}

\noindent\textbf{Keywords:} Rainbow; Rainbow~2; hash functions; collision attacks; chosen-prefix collisions; non-cryptographic hashing; state collapse; SMHasher3; AI-assisted cryptanalysis.

\section{Introduction}
\label{sec:introduction}

Rainbow is a word-oriented hash with a 256-bit internal state, 16-byte input blocks, a 64-bit seed, and 64-, 128-, or 256-bit digests. Its documentation is explicit that it is a general-purpose, non-cryptographic hash and is ``not designed as cryptographic''~\cite{rain-readme}. That disclaimer is important and we take it seriously: nothing in this paper contradicts a claim Rainbow makes.

It is worth stating plainly why a collision attack on such a function is of interest at all. Collisions in fast non-cryptographic hashes are usually easy, and often far easier than the attack presented here: CRC32 is linear and collides by construction, and published collision constructions exist for xxHash, MurmurHash3, and CityHash at negligible cost. A $2^{32}$ search against a 64-bit digest is simply the generic birthday bound and is not a cryptanalytic result.

The result here is of a different kind, and rests on two observations.

First, the attack does not target a digest; it targets the internal state, and it succeeds against the \emph{full} function. Rainbow maintains 256 bits of state, but we show that only 128 of those bits can ever distinguish two states from the standpoint of an attacker who controls one more block, because message injection preserves two 64-bit quantities and lets the attacker place the state anywhere else in their coset. The state does not merely mix poorly; it is, for collision purposes, 128 bits wide.

Second, those 128 bits split. The $\mixA$ mixer never moves information between $(h_0,h_1)$ and $(h_2,h_3)$, so the two pair sums are refreshed by two separate 64-bit maps. Divide and conquer applies, and the cost of merging two arbitrary states is two 64-bit searches -- roughly $2^{33}$ -- instead of the $2^{128}$ that a 256-bit state should demand.

The consequence is that the digest width is decorative with respect to collisions. The same pair of messages collides at 64, 128, and 256 bits, because the states are already equal before finalization begins. Rainbow-256 costs no more to collide than rainbow-64.

\paragraph{Contributions.}
\begin{itemize}
  \item A proof that two Rainbow states merge under one message block, with zero search, exactly when their two pair sums agree (Theorem~\ref{thm:pairsum}), validated on $1000/1000$ random state pairs.
  \item A reduction of full-state collision search to two independent 64-bit problems (Section~\ref{sec:independence}), yielding a same-seed collision in $2^{32.28}$ evaluations (Section~\ref{sec:sameiv}) and a chosen-prefix collision in $2^{34.80}$ evaluations (Section~\ref{sec:chosenprefix}), verified across four independent implementations (Section~\ref{sec:verification}).
  \item Characterization of three further structural weaknesses in Rainbow v3.7.1 (Section~\ref{sec:further}): the 192-bit finalization bottleneck in which word $h_0''$ is $100\%$ dead in 64-bit mode and the $w=4$ toy finalizer reaches only $6.1\%$ of the state space (\code{Q20}), 2-adic differential filtration rays (\code{Q18}--\code{Q19}), and top-bit seed IV linearity (\code{Q22}).
  \item A seven-rung hardening ladder (\code{R0}--\code{R6}) evaluated algebraically and empirically through 32-core cloud \code{SMHasher3} campaigns on AMD EPYC Milan (Section~\ref{sec:ladder}), including the proof and elimination of a bit-0 parity cancellation ($(Y \oplus X) - X \equiv Y \pmod 2$) between post-multiply Feistel cross-XORs and unrotated four-lane squeezes (Proposition~\ref{prop:squeezecancel}), culminating in the selection of \code{R5} as \textbf{Rainbow~2} (\code{rainbow2}).
\end{itemize}

\section{Background: Rainbow v3.7.1}
\label{sec:background}

We describe the function as implemented in \code{src/rainbow.cpp} at version 3.7.1, which is the authority for every claim below. All arithmetic is in $\W$; $\ROTR(x,n)$ is a 64-bit rotation.

\subsection{Constants and mixers}

Rainbow uses eight odd 64-bit constants $P,Q,R,S,T,U,V,W$, with $P = 2^{64}-59$. The state is $h = (h_0,h_1,h_2,h_3)$.

The primary mixer $\mixA$ is
\begin{align}
  a &\leftarrow \ROTR(h_0 \cdot P,\ 23)\cdot Q, \label{eq:mixa-a}\\
  b &\leftarrow \ROTR\big((h_1 \oplus a) \cdot R,\ 29\big)\cdot S, \label{eq:mixa-b}\\
  c &\leftarrow \ROTR(h_2 \cdot T,\ 31)\cdot U, \label{eq:mixa-c}\\
  d &\leftarrow \ROTR\big((h_3 \oplus c) \cdot V,\ 37\big)\cdot W, \label{eq:mixa-d}
\end{align}
with $h \leftarrow (a,b,c,d)$. The secondary mixer $\mixB$, which takes the seed as an additional input, rewrites only $h_1$ and $h_2$ and is followed by a rotation of the four words.

\begin{remark}[$\mixA$ is pairwise separable]
\label{rem:separable}
Equations~\eqref{eq:mixa-a}--\eqref{eq:mixa-b} read only $h_0,h_1$ and write only $a,b$; equations~\eqref{eq:mixa-c}--\eqref{eq:mixa-d} read only $h_2,h_3$ and write only $c,d$. $\mixA$ therefore acts as two independent maps on $\W^2$. This is the property the attack exploits, and it is visible directly in the source.
\end{remark}

\subsection{Message schedule}

For each complete 16-byte block, Rainbow parses two little-endian words $x,y$ and injects
\begin{equation}
  h_0 \mathrel{-}= x,\qquad h_1 \mathrel{+}= x,\qquad h_2 \mathrel{+}= y,\qquad h_3 \mathrel{-}= y.
  \label{eq:inject}
\end{equation}
Blocks alternate mixers: the first complete block is followed by $\mixA$, the second by $\mixB$ and a rotation, and so on. Writing blocks $1$-indexed, block $n$ is followed by $\mixA$ exactly when $n$ is odd.

The state is keyed by the total message length: $h_i = \mathrm{seed} + \mathrm{len} + \{1,2,3,5\}$ for $i=0,1,2,3$. Two messages therefore share an initial state if and only if they share a seed and a length, and every collision below is between equal-length messages.

After the block loop, Rainbow applies $\mixB$, absorbs the $0$--$15$ tail bytes at length-determined lanes, then applies $\mixA$, $\mixB$, $\mixA$ and emits digest words.

\begin{remark}[The output is itself a pair sum]
\label{rem:output}
Each 64-bit digest word is computed as \code{g = 0; g -= h[2]; g -= h[3];}, that is $-\sigma_2$ for the current state. The finalization reads out precisely the quantity that the message schedule preserves. We do not need this observation for the attack, which achieves full-state equality, but it indicates that the pair sums are load-bearing at both ends of the construction.
\end{remark}

\section{Threat Model, Notation, and Scope}
\label{sec:scope}

We work in the fixed public-seed setting with $\mathrm{seed}=0$; nothing depends on that choice, and the attack is stated for arbitrary known seeds. The attacker chooses complete messages and knows the seed. There is no key and no secret, so no PRF or MAC claim is in scope.

Write $\sigma_1(h) = h_0 + h_1$ and $\sigma_2(h) = h_2 + h_3$ for the two \emph{pair sums}. We say two states are \emph{mergeable} if some choice of one further message block for each yields identical states.

We attack the full, unreduced function: every mixer, the complete schedule, the finalization, and all three digest widths. No component is weakened, removed, or reduced in round count, and the attack code does not link against production sources, so it cannot perturb them.

\section{The Pair-Sum Invariant}
\label{sec:invariant}

\begin{lemma}[Injection preserves both pair sums]
\label{lem:preserve}
Let $h'$ be the result of injecting $(x,y)$ into $h$ by \eqref{eq:inject}. Then $\sigma_1(h') = \sigma_1(h)$ and $\sigma_2(h') = \sigma_2(h)$, for every $x,y \in \W$.
\end{lemma}

\begin{proof}
$\sigma_1(h') = (h_0 - x) + (h_1 + x) = h_0 + h_1$ and $\sigma_2(h') = (h_2 + y) + (h_3 - y) = h_2 + h_3$, all in $\W$.
\end{proof}

The message words move the state along the two cosets and cannot leave them. The attacker's freedom inside a coset is total: $x$ and $y$ range over all of $\W$. This makes the pair sums an exact invariant of the schedule and, as the next theorem shows, a complete one.

\begin{theorem}[Pair-sum equivalence]
\label{thm:pairsum}
Let $S$ and $T$ be Rainbow states at the same block index. There exist message words $(x,y)$ for $S$ and $(x',y')$ for $T$ whose injection yields identical states if and only if
\[
  \sigma_1(S) = \sigma_1(T) \quad\text{and}\quad \sigma_2(S) = \sigma_2(T).
\]
When the condition holds, the solution set has size $2^{128}$: $x$ and $y$ may be chosen freely, and then
\[
  x' = x + (t_0 - s_0), \qquad y' = y + (t_3 - s_3).
\]
\end{theorem}

\begin{proof}
($\Leftarrow$) Injection yields $(s_0-x,\, s_1+x,\, s_2+y,\, s_3-y)$ and $(t_0-x',\, t_1+x',\, t_2+y',\, t_3-y')$. Matching word $0$ requires $x' = x + (t_0-s_0)$; matching word $1$ requires $x' = x + (s_1-t_1)$. The two are consistent precisely when $t_0-s_0 = s_1-t_1$, i.e. $s_0+s_1 = t_0+t_1$. Symmetrically, word $3$ requires $y' = y + (t_3-s_3)$ and word $2$ requires $y' = y + (s_2-t_2)$, consistent precisely when $s_2+s_3 = t_2+t_3$. Under the hypothesis both hold, and $x,y$ remain free.

($\Rightarrow$) If the post-injection states are identical they have equal pair sums, and by Lemma~\ref{lem:preserve} injection did not change either state's pair sums, so $S$ and $T$ had equal pair sums already.
\end{proof}

\begin{corollary}[Effective state width]
\label{cor:width}
For collision finding, the collision-relevant content of a Rainbow state is the 128-bit pair $(\sigma_1,\sigma_2)$, not the 256-bit state. Two states agreeing on it are one attacker-chosen block away from being the same state, at no search cost.
\end{corollary}

Note the direction of the ``if and only if'': the theorem is not merely a sufficient trick but a characterization. The residual $128$ bits of state beyond the pair sums are attacker-erasable.

\begin{result}[Theorem validation]
\label{res:theorem}
Over $1000$ random trials, state pairs were constructed with independently random words subject only to matching pair sums, and merged by the closed form above. All $1000$ produced bitwise-identical states after injection, and identical states after each of $\mixA$ and $\mixB$ with rotation. Implemented in \code{research/rainbow/chosen\_prefix\_collision.cpp} as \code{theorem\_selftest}.
\end{result}

\section{Independence of the Two Halves}
\label{sec:independence}

Corollary~\ref{cor:width} reduces the merging problem to matching $128$ bits. Remark~\ref{rem:separable} reduces it further, to two independent $64$-bit problems.

\begin{definition}[Half-sum maps]
For a state $S$ at a $\mixA$ boundary, define
\begin{align*}
  A_S(x) &= \ROTR\big((s_0-x)\cdot P,\,23\big)\cdot Q, &
  B_S(x) &= \ROTR\big(((s_1+x)\oplus A_S(x))\cdot R,\,29\big)\cdot S,\\
  C_S(y) &= \ROTR\big((s_2+y)\cdot T,\,31\big)\cdot U, &
  D_S(y) &= \ROTR\big(((s_3-y)\oplus C_S(y))\cdot V,\,37\big)\cdot W,
\end{align*}
and set $H_S(x) = A_S(x) + B_S(x)$ and $K_S(y) = C_S(y) + D_S(y)$.
\end{definition}

$H_S$ is the value of $\sigma_1$ after injecting $x$ and applying $\mixA$; $K_S$ is the value of $\sigma_2$ after injecting $y$ and applying $\mixA$. By Remark~\ref{rem:separable}, $H_S$ depends only on $x$ and $K_S$ only on $y$.

\begin{proposition}[Two-claw merging]
\label{prop:twoclaw}
Let $S,T$ be states at a $\mixA$ boundary. If $H_S(x) = H_T(x')$ and $K_S(y) = K_T(y')$, then injecting $(x,y)$ into $S$ and $(x',y')$ into $T$, applying $\mixA$ to each, and injecting the bridge words
\[
  p = -B_S(x),\quad q = -C_S(y) \qquad\text{and}\qquad p' = -B_T(x'),\quad q' = -C_T(y')
\]
produces identical states, namely $(H,\,0,\,0,\,K)$.
\end{proposition}

\begin{proof}
After $\mixA$ the first pair is $(A,B)$ and the second is $(C,D)$. The bridge injection sends $(A,B)\mapsto(A-p,\,B+p)$ and $(C,D)\mapsto(C+q,\,D-q)$. With $p=-B$ the first pair becomes $(A+B,0)=(H,0)$; with $q=-C$ the second becomes $(0,C+D)=(0,K)$. The hypotheses make the two $H$ values and the two $K$ values equal, so the four words agree. (Any other common point of the cosets would serve equally; this one is convenient because it is canonical.)
\end{proof}

Since $H$ and $K$ behave as random $64$-bit maps, each condition is met after roughly $2^{32}$ evaluations by birthday search, for a total near $2^{33}$ rather than $2^{128}$.

\begin{remark}[$H$ is not a permutation]
$A_S$ is a bijection of $x$, and $t\mapsto\ROTR(tR,29)S$ is a bijection, so one might hope $H_S = A_S + B_S$ is a permutation, which would make the single-search variant of Section~\ref{sec:sameiv} impossible. It is not: the collision exhibited in Result~\ref{res:sameiv} is an explicit pair $x_a \neq x_b$ with $H(x_a)=H(x_b)$, and the search behaved throughout like a search on a random map. The XOR in \eqref{eq:mixa-b} sits between two additive structures and destroys any such algebraic balance.
\end{remark}

\section{Same-Seed Collision on Full Rainbow}
\label{sec:sameiv}

When the two messages share a seed and a length, they share an initial state, and one search suffices.

\paragraph{Construction.} Fix message length $32$ bytes (two blocks) and seed $0$, giving $S=(33,34,35,37)$. Let both messages use the same second word $y=0$ in block~1, so the second pair evolves identically and $K$ agrees automatically. Search for $x_a \neq x_b$ with
\[
  H(x_a) = H(x_b).
\]
Emit
\[
  m_i = \mathrm{LE}_{64}(x_i)\ \|\ \mathrm{LE}_{64}(0)\ \|\ \mathrm{LE}_{64}(-B(x_i))\ \|\ \mathrm{LE}_{64}(0).
\]
Block~2's first word cancels $B$, so both states become $(H,0,C,D)$ with $C,D$ common. The states are equal before block~2's mixer, hence for the rest of the computation.

\paragraph{Search.} We use a van Oorschot--Wiener parallel distinguished-point rho on $H$~\cite{vow1999}: walks iterate $x \mapsto H(x)$ until the low $d$ bits vanish, endpoints are stored, and a repeated endpoint is re-walked to expose the colliding pair. Memory is negligible.

\begin{result}[Same-seed collision]
\label{res:sameiv}
With $d=20$ and 8 threads, a collision was found in $5{,}206{,}393{,}642$ evaluations ($2^{32.28}$), $4920$ walks, one endpoint merge, $6.63$\,s wall, $785.7$ Meval/s.
\begin{center}\footnotesize
\begin{tabularx}{\linewidth}{@{}lX@{}}
\toprule
$x_a$ & \code{0xf2361ca7274d362f} \\
$x_b$ & \code{0x17f81881806268c1} \\
$H$   & \code{0x26d62ded70d34542} (both) \\
$m_a$ & \code{2f364d27a71c36f20000000000000000dbc4cde348f671ec0000000000000000} \\
$m_b$ & \code{c16862808118f8170000000000000000ebc98b7034e669320000000000000000} \\
\midrule
rainbow-64  & \code{b013b427e381f006} \\
rainbow-128 & \code{b013b427e381f006c89c06357350be4b} \\
rainbow-256 & \code{b013b427e381f006c89c06357350be4be142f3bdc749f11c58acc1b99a4d91e3} \\
\bottomrule
\end{tabularx}
\end{center}
The merged state after block-2 injection is $(\code{26d62ded70d34542},\,0,\,\code{1b0560be4a530f25},\,\code{b8edca491361c74b})$ for both messages.
\end{result}

\section{Chosen-Prefix Collision}
\label{sec:chosenprefix}

The same-seed attack requires the messages to agree before the collision block. Theorem~\ref{thm:pairsum} and Proposition~\ref{prop:twoclaw} remove that restriction entirely: \emph{any} two states can be merged, so the prefixes may be chosen arbitrarily and independently.

\paragraph{Construction.} Let $\mathrm{pre}_a \neq \mathrm{pre}_b$ be prefixes of equal length, padded to a multiple of $32$ bytes so that the next block is a $\mixA$ block. Absorb them to obtain $S$ and $T$. Solve two claws,
\[
  H_S(x_a) = H_T(x_b), \qquad K_S(y_a) = K_T(y_b),
\]
then append to each message its claw block $\mathrm{LE}_{64}(x)\|\mathrm{LE}_{64}(y)$ and its bridge block $\mathrm{LE}_{64}(-B)\|\mathrm{LE}_{64}(-C)$. By Proposition~\ref{prop:twoclaw} the states coincide at the bridge injection.

A claw between two functions is obtained as an ordinary collision of the flavored combination $F(w) = H_{\beta(w)}(w)$, where $\beta(w)$ is a cheap bit derived from $w$; collisions with differing flavor are claws. Roughly half the collisions are same-flavor and discarded, so a claw costs about twice a plain collision. This accounts for the observed $2^{33.9}$ per claw against the $2^{32}$ of Result~\ref{res:sameiv}; a dedicated two-sided search would reduce it.

\begin{result}[Chosen-prefix collision]
\label{res:chosenprefix}
Two 96-byte messages with opposed meanings and identical digests at all three widths. Prefixes were space-padded to 64 bytes; the padding is part of the chosen prefix.
\begin{center}\footnotesize
\begin{tabularx}{\linewidth}{@{}lX@{}}
\toprule
prefix $a$ & \code{From: alice@example.com\ \ Amount: USD\ \ \ \ \ \ 10.00} \\
prefix $b$ & \code{From: mallory@example.net Amount: USD 1000000.00} \\
\midrule
claw $H$ & $16{,}190{,}842{,}307$ evaluations ($2^{33.91}$), $27.1$\,s \\
claw $K$ & $13{,}673{,}595{,}949$ evaluations ($2^{33.67}$), $23.0$\,s \\
total    & $29{,}864{,}438{,}256$ evaluations ($2^{34.80}$), $50.0$\,s \\
\midrule
merged state & $(\code{0bcacc509d3457d6},\,0,\,0,\,\code{7f92191466aadf76})$ \\
rainbow-64  & \code{913ac822c5f4c8fe} \\
rainbow-128 & \code{913ac822c5f4c8fe40cb08c0af232d8a} \\
rainbow-256 & \code{913ac822c5f4c8fe40cb08c0af232d8ab58e007cc27e827269d591fd423682db} \\
\bottomrule
\end{tabularx}
\end{center}
\end{result}

\subsection{Corollaries}

Because a merge produces genuinely identical internal states, the standard generic consequences follow immediately and at the same per-merge cost.

\begin{itemize}
  \item \textbf{Multicollisions.} Chaining $t$ successive same-state collision gadgets in the manner of Joux~\cite{joux2004} yields $2^t$ colliding messages for $t \cdot 2^{32}$ work. A $2^{64}$-multicollision costs about $2^{38}$.
  \item \textbf{Common-suffix forgery.} Any common suffix may be appended to a merged pair, at every digest width, since the states are equal.
  \item \textbf{Herding and expandable messages.} The cheap merge of arbitrary states is the primitive underlying diamond-structure herding~\cite{kelseykohno2006} and long-message second-preimage techniques~\cite{kelseyschneier2005}; we did not implement these.
\end{itemize}

We emphasize that only the first two items were executed; the third is stated as an expected consequence, not a demonstrated result.

\section{Why the Digest Width Does Not Help}
\label{sec:width}

Rainbow's finalization is a deterministic function of the internal state. The attack achieves equality of that state before finalization begins, so every subsequent operation, and hence every digest word at every width, agrees. This is visible in the witnesses: the 64-bit digest is a prefix of the 128-bit digest, which is a prefix of the 256-bit digest, for both messages.

\begin{center}
\begin{tabular}{@{}lccc@{}}
\toprule
Variant & Generic collision cost & Cost here & Shortfall \\
\midrule
rainbow-64  & $2^{32}$  & $2^{32.3}$ & none (generic) \\
rainbow-128 & $2^{64}$  & $2^{32.3}$ & $\approx 2^{32}$ \\
rainbow-256 & $2^{128}$ & $2^{32.3}$ & $\approx 2^{96}$ \\
\bottomrule
\end{tabular}
\end{center}

The honest reading is narrow and, we think, still worth stating. Rainbow-64 is not weakened: $2^{32}$ is what a 64-bit digest costs. What fails is the implicit proposition that selecting a wider digest buys proportionally more collision resistance. It buys none. A user choosing rainbow-256 over rainbow-64 for collision avoidance in a large data set receives no benefit against an adversary, and the cost of demonstrating this is seconds on a laptop.

This is a property of the construction, not of the search: a 256-bit state whose message schedule preserves two 64-bit invariants, and whose mixer never crosses the halves, cannot offer more than $64$-bit-per-half resistance no matter how many finalization rounds follow.

\section{Experimental Results and Independent Verification}
\label{sec:verification}

Every witness was checked by four implementations that do not share code:

\begin{enumerate}
  \item \textbf{The attack's own model.} \code{collision\_pair\_sum.cpp} and \code{chosen\_prefix\_collision.cpp} each transcribe Rainbow independently and self-check their witnesses.
  \item \textbf{An independent Python implementation.} \code{verify\_pair\_sum.py} re-implements Rainbow with big integers, re-derives the pair sums from the message bytes alone, prints a checkpoint-by-checkpoint comparison, and confirms the digests. It reports the merge point: block-2 injection for Result~\ref{res:sameiv}, block-6 injection for Result~\ref{res:chosenprefix}.
  \item \textbf{The unmodified production source.} \code{replay\_production.cpp} includes \code{src/rainbow.cpp} verbatim and hashes the serialized bytes, never in-memory words.
  \item \textbf{The shipped command-line tool.} \code{rainsum -a bow -s \{64,128,256\}} on the message files reproduces the same digests.
\end{enumerate}

All four agree on both witnesses at all three widths. Messages are carried as hex and replayed from bytes, so no endianness assumption is shared between the search and the verification.

\section{Further Structural Weaknesses in Rainbow v3.7.1}
\label{sec:further}

Beyond the pair-sum coset invariant (\code{W1}) and the pair-separability of $\mixA$ (\code{W2}), our algebraic spectrum campaign (\code{research/rainbow-spectrum/}, \code{Q01}--\code{Q28}) identified three additional structural weaknesses in Rainbow v3.7.1 (\code{W3}--\code{W5}) that must be accounted for in any hardened successor.

\subsection{W3: The 192-Bit Finalization Bottleneck (\code{Q20})}
\label{sec:bottleneck}

After all 16-byte blocks are processed, Rainbow v3.7.1 executes $\mixB(h)$, absorbs any $0 \le r < 16$ tail bytes into $h$, and runs the finalization sequence
\[
  h^{(1)} = \mixA(h), \qquad h'' = \mixB(h^{(1)}, \text{seed}), \qquad h^{(3)} = \mixA(h''),
\]
followed by the output squeeze $g_0 = 0 - h^{(3)}_2 - h^{(3)}_3$ for 64-bit digests, and an additional $h^{(4)} = \mixA(h^{(3)})$ with $g_1 = 0 - h^{(4)}_2 - h^{(4)}_3$ for 128-bit digests (and a second $\mixA \circ \mixB \circ \mixA$ cycle for words $g_2, g_3$ of 256-bit digests).

\begin{proposition}[Finalization dead-word fibre]
\label{prop:deadword}
In Rainbow v3.7.1:
\begin{enumerate}
  \item Let $h'' = (h''_0, h''_1, h''_2, h''_3)$ be the internal state immediately after $\mixB$ in finalization. Because $\mixA$ is pair-separable and the squeeze $g = -h_2 - h_3$ reads only lanes $(2,3)$, neither $g_0$ (the entire 64-bit digest) nor $g_1$ (the second 64-bit word of the 128-bit digest) depends on $(h''_0, h''_1)$.
  \item Moreover, because $\mixB$ leaves lane $0$ untouched ($h''_0 = h^{(1)}_0$) and mixes only lanes $(1,2)$ via $a = \ROTR(h^{(1)}_1 V, 29)W$ and $b = \ROTR((h^{(1)}_2 \oplus (a + \text{seed}))R, 29)S$, the post-$\mixA$ word $h^{(1)}_0$ is $100\%$ dead in 64-bit and 128-bit finalization: perturbing $h^{(1)}_0$ arbitrarily after the first finalization $\mixA$ changes $0/2000$ digests in \code{rainbow-64} and \code{rainbow-128}, and perturbing $h_0''$ before the middle $\mixB$ in \code{rainbow-256} changes $0/2000$ of the upper digest words $(g_2, g_3)$.
\end{enumerate}
\end{proposition}

\begin{proof}
Write $\mixA(s_0, s_1, s_2, s_3) = (F_1(s_0, s_1), F_2(s_2, s_3))$. After $h'' = \mixB(h^{(1)}, \text{seed})$, we have $h^{(3)} = (F_1(h''_0, h''_1), F_2(h''_2, h''_3))$ and $h^{(4)} = (F_1^2(h''_0, h''_1), F_2^2(h''_2, h''_3))$. Since $g_0 = -(h^{(3)}_2 + h^{(3)}_3)$ and $g_1 = -(h^{(4)}_2 + h^{(4)}_3)$ depend solely on $F_2(h''_2, h''_3)$ and $F_2^2(h''_2, h''_3)$, the 128-bit half $(h''_0, h''_1)$ never enters $g_0$ or $g_1$. Because $\mixB$ maps $(h^{(1)}_0, h^{(1)}_1, h^{(1)}_2, h^{(1)}_3) \mapsto (h^{(1)}_0, b, a, h^{(1)}_3)$ where $a,b$ depend only on $(h^{(1)}_1, h^{(1)}_2)$, $h^{(1)}_0$ passes into $h''_0$ without affecting $(h''_2, h''_3)$, compressing the 256-bit state through a strict 192-bit bottleneck $(h^{(1)}_1, h^{(1)}_2, h^{(1)}_3)$.
\end{proof}

On the scaled $w=4$ toy model (4-bit lanes, 16-bit total state space of $65{,}536$ states), exhaustive enumeration over all $65{,}536$ pre-finalization states confirms that Rainbow v3.7.1's finalizer reaches only $4{,}007 / 65{,}536$ ($6.11\%$) distinct 4-word digest vectors, compared to $1 - e^{-1} \approx 63.21\%$ for a uniform random map.

\subsection{W1b: 2-Adic Differential Filtration and DDT Ray Anisotropy (\code{Q18}--\code{Q19})}
\label{sec:ddt}

Every constant multiplication $x \mapsto x \cdot C \bmod 2^w$ by an odd constant $C$ is an automorphism of the 2-adic filtration: if $x \equiv x' \pmod{2^k}$, then $x C \equiv x' C \pmod{2^k}$, so any difference $\Delta = 2^k u$ ($u$ odd) preserves its 2-adic valuation $\nu_2(\Delta) = k$ through multiplication. Only the fixed rotations $\ROTR(\cdot, r)$ move bits across valuation levels. In Rainbow v3.7.1, the symmetric block injection $\Delta h = (-\Delta g_0 - \Delta g_1,\, \Delta g_0 + \Delta g_1,\, \Delta g_1,\, -\Delta g_1)$ feeds aligned high-valuation differences simultaneously into both pairs, producing sharp anisotropic rays in the 1-block differential distribution table (DDT). At $w=12$ (exhaustive over $2^{12}$ block differences), the maximum 2-adic ray collision factor reaches $M_r = 11.01\times$ above uniform and the quotiented ray concentration reaches $S_r = 145.29\times$.

\subsection{W4 \& W5: Cross-Width Prefixing and Top-Bit Seed Linearity (\code{Q22})}
\label{sec:seediv}

Rainbow v3.7.1 initializes its four state lanes identically across all digest widths as
\[
  h_0 = \text{seed} + \text{olen} + 1,\quad h_1 = \text{seed} + \text{olen} + 2,\quad h_2 = \text{seed} + \text{olen} + 3,\quad h_3 = \text{seed} + \text{olen} + 5.
\]
Two structural consequences follow:
\begin{enumerate}
  \item \textbf{Cross-width output prefixing (\code{W4}):} Because \code{hashsize} is absent from the IV and finalization emits words sequentially from the same state trajectory, \code{rainbow-64}$(m)$ is always the exact 8-byte prefix of \code{rainbow-128}$(m)$ and \code{rainbow-256}$(m)$, and \code{rainbow-128}$(m)$ is the exact 16-byte prefix of \code{rainbow-256}$(m)$ ($200/200$ trials).
  \item \textbf{Top-bit related-seed linearity (\code{W5}, \code{Q22}):} Flipping the most significant bit of the seed ($\Delta\text{seed} = 2^{63}$) flips only bit 63 in each initial lane $h_i$ (since $+2^{63} \equiv \oplus 2^{63} \pmod{2^{64}}$ generates no carry above bit 63), giving an initial 256-bit state Hamming weight of exactly $4.0$ bits. Moreover, when $h_i$ is first multiplied by an odd prime ($P$ or $T$ in $\mixA$, $V$ in $\mixB$), $(h_i \oplus 2^{63}) \cdot P = h_i P + 2^{63} P \equiv h_i P \oplus 2^{63} \pmod{2^{64}}$ because $P$ is odd, so the single-bit difference survives the first multiplication with probability $1$ before rotation! At $w=8$, $9.5\%$ of random message collisions transfer intact from $\text{seed}$ to $\text{seed} \oplus 2^{w-1}$ ($19\times$ above random chance).
\end{enumerate}

\section{The Hardening Ladder (\code{R0}--\code{R6}) and the Design of Rainbow~2}
\label{sec:ladder}

Rather than proposing an untested patch or sacrificing Rainbow's speed with heavy extra rounds, we constructed a seven-rung hardening ladder (\code{R0}--\code{R6}) in \code{research/rainbow2/} to isolate the exact algebraic and cycle-count effect of each structural modification.

\subsection{Specification of the Seven Hardening Rungs}
\label{sec:rungs}

\begin{itemize}
  \item \textbf{\code{R0} (Baseline Rainbow v3.7.1):} Unmodified \code{src/rainbow.cpp} control.
  \item \textbf{\code{R1} (Asymmetric Rotated Block Injection --- fixes \code{W1}, \code{W1b}):} Replaces the rank-2 pair-sum block load with the full-rank asymmetric load
  \[
    h_0 \mathrel{-}= g_0,\qquad h_1 \mathrel{+}= \ROTR(g_0, 19),\qquad h_2 \mathrel{+}= g_1,\qquad h_3 \mathrel{-}= \ROTR(g_1, 41).
  \]
  Because $x \mapsto (-x, \ROTR(x, 19))$ has odd determinant $-(1 - \ROTR(\cdot, 19))$ modulo 2 on non-constant bits, the map $(g_0, g_1) \mapsto (\Delta\sigma_1, \Delta\sigma_2)$ is a bijection over $(\W)^2$. One-block state merges drop from $1000/1000$ to $0/1000$, and the $w=12$ DDT ray bias $M_r$ drops from $11.01\times$ to $3.95\times$.
  \item \textbf{\code{R2} (Feistel Cross-Coupled $\mixA$ --- fixes \code{W2}, \code{W3}):} Appends the two reversible cross-half XORs
  \[
    a \mathrel{\oplus}= d,\qquad c \mathrel{\oplus}= b
  \]
  at the end of $\mixA$ before writing $(s_0, s_1, s_2, s_3) = (a, b, c, d)$. This preserves exact bijectivity of $\mixA$ while achieving $0.500$ cross-pair avalanche after two $\mixA$ calls, making $h^{(1)}_0$ and $(h''_0, h''_1)$ $100\%$ live in finalization ($2000/2000$ trials) and expanding the $w=4$ finalizer image from $4{,}007$ ($6.1\%$) to $41{,}766$ ($63.7\%$).
  \item \textbf{\code{R3} (\code{R1 + R2} Combined Core):} Combines \code{R1}'s asymmetric rotated block injection with \code{R2}'s cross-coupled $\mixA$.
  \item \textbf{\code{R4} (\code{R3} + Rotated Four-Lane Output Squeeze):} Replaces the two-lane squeeze $g = 0 - h_2 - h_3$ with the four-lane rotated squeeze
  \[
    g = \ROTR(h_0 - h_1, 29) - h_2 - h_3,
  \]
  so every output word directly combines all four 64-bit lanes.
  \item \textbf{\code{R5} (\code{R4} + Nonlinear Domain-Separated Rotated IV --- fixes \code{W4}, \code{W5}; selected as \textbf{Rainbow~2}):} Replaces the linear additive IV with a domain-separated, pre-rotated prime-multiply initialization binding \code{seed}, \code{olen}, and \code{hashsize}:
  \begin{align*}
    \mathrm{dom} &= (\mathrm{uint64\_t}(\mathrm{hashsize}) \ll 32) \oplus (\mathrm{olen} \cdot P), \\
    h_0 &= \ROTR(\mathrm{seed} \oplus \mathrm{dom}, 17) \cdot Q + 1, &
    h_1 &= (\ROTR(\mathrm{seed} + \mathrm{olen}, 29) \cdot R) \oplus (\mathrm{dom} + 2), \\
    h_2 &= \ROTR(\mathrm{seed} \oplus (\mathrm{olen} \cdot T), 41) \cdot U + 3, &
    h_3 &= (\ROTR(\mathrm{seed} + \mathrm{dom}, 53) \cdot V) \oplus W.
  \end{align*}
  Rotating \code{seed} \emph{before} multiplying by $\{Q,R,U,V\}$ moves bit 63 into interior bit positions $\{46, 34, 22, 10\}$ prior to multiplication, increasing the $\Delta\text{seed}=2^{63}$ IV Hamming weight from $4.0$ bits to $76.7$ bits and reducing $w=8$ related-seed collision transfer from $9.5\%$ to $0.5\%$ (random chance), while eliminating cross-width output prefixing ($0/200$).
  \item \textbf{\code{R6} (Zero-Instruction Cyclic Cross-Coupled $\mixA$):} Moves the cross-XORs \emph{before} the second multiply chain in $\mixA$ ($b \mathrel{\oplus}= c;\ d \mathrel{\oplus}= a;$ prior to $b \cdot R \dots S$ and $d \cdot V \dots W$) and applies a compile-time cyclic output rename $(s_0, s_1, s_2, s_3) = (b, c, d, a)$, paired with \code{R1} injection and $g = (h_0 \oplus h_2) - (h_1 \oplus h_3)$.
\end{itemize}

\subsection{The Iterative Algebra $\leftrightarrow$ SMHasher3 Loop: Parity Cancellation in Unrotated Squeezes}
\label{sec:squeezecancel}

Our initial specification of \code{R4} and \code{R5} paired \code{R2}'s post-multiply Feistel $\mixA$ ($h_0 = a \oplus d,\ h_1 = b,\ h_2 = c \oplus b,\ h_3 = d$) with the unrotated XOR-subtract squeeze $g_{\text{unrot}} = (h_0 \oplus h_2) - (h_1 \oplus h_3)$. Although $g_{\text{unrot}}$ passed all $w=4$ and $w=64$ algebraic rank and dead-word tests, running the 32-core \code{SMHasher3} \code{Avalanche} and \code{BIC} suites immediately flagged a severe bias ($27\%$--$77\%$) at output bit~0 (\code{out 0}) on short non-block-aligned keys ($3$--$12$ and $20$ bytes), whereas \code{R3} and \code{R6} passed cleanly ($<0.95\%$). Inspecting the interaction between $\mixA$ and $g_{\text{unrot}}$ reveals an exact algebraic identity:

\begin{proposition}[Bit-0 parity cancellation in unrotated Feistel squeezes]
\label{prop:squeezecancel}
Let $(a, b, c, d)$ be the four lane registers after the multiplicative chains of $\mixA$ in \code{R2}, so the output of $\mixA$ is $(h_0, h_1, h_2, h_3) = (a \oplus d,\, b,\, c \oplus b,\, d)$. Set $Y = a \oplus c$ and $X = b \oplus d$. Then the unrotated squeeze $g_{\text{unrot}} = (h_0 \oplus h_2) - (h_1 \oplus h_3)$ satisfies
\[
  g_{\text{unrot}} = (Y \oplus X) - X \pmod{2^{64}},
\]
and its least significant bit is identically
\[
  (g_{\text{unrot}})_0 \equiv (Y_0 \oplus X_0) \oplus X_0 \equiv Y_0 \equiv a_0 \oplus c_0 \pmod 2,
\]
which is \textbf{completely independent of $b$ and $d$}.
\end{proposition}

\begin{proof}
Substituting $h_0 = a \oplus d$ and $h_2 = c \oplus b$ gives $h_0 \oplus h_2 = (a \oplus c) \oplus (b \oplus d) = Y \oplus X$, while $h_1 \oplus h_3 = b \oplus d = X$. Modulo $2$, integer subtraction coincides with XOR, so the bit-0 contribution of $X = b \oplus d$ cancels identically. In Rainbow's tail switch (\code{len} $\in \{1,\dots,15\}$), bytes injected into $h_1$ or $h_3$ enter the final $\mixA$ solely through $b$ and $d$; because $a$ and $c$ do not depend on $b$ or $d$ prior to the final $a \mathrel{\oplus}= d,\ c \mathrel{\oplus}= b$ step, canceling $X = b \oplus d$ in bit 0 erases those tail bytes from $(g_{\text{unrot}})_0$ unless a borrow propagates from lower bits.
\end{proof}

Replacing $g_{\text{unrot}}$ in \code{R4} and \code{R5} with the rotated four-lane squeeze $g = \ROTR(h_0 - h_1, 29) - h_2 - h_3$ misaligns the bit positions of $(h_0, h_1)$ relative to $(h_2, h_3)$ by 29 bits and replaces the shared XOR term with modular subtraction, eliminating the cancellation completely and reducing maximum \code{SMHasher3} \code{Avalanche} bias on \code{R4} and \code{R5} from $77.27\%$ down to $0.68\%$--$0.92\%$ across all key lengths and widths.

\subsection{Algebraic and Empirical SMHasher3 Evaluation}
\label{sec:results-ladder}

Tables~\ref{tab:algebraic-ladder} and~\ref{tab:smhasher3-ladder} summarize the complete algebraic, cryptanalytic, and 32-vCPU AMD EPYC Milan (\code{c2d-highcpu-32}) \code{SMHasher3} results across all seven rungs (\code{R0}--\code{R6}) and all three digest widths ($64, 128, 256$ bits). Every variant was first verified for $100\%$ streaming vs.\ single-call equivalence across $15{,}960$ chunk/length/seed combinations under \code{AddressSanitizer} and \code{UndefinedBehaviorSanitizer}, and $2{,}352$ native C++ vs.\ pure-Python reference checks.

\begin{table}[ht!]
\centering
\caption{Algebraic and cryptanalytic verification across the seven hardening rungs (\code{R0}--\code{R6}).}
\label{tab:algebraic-ladder}
\vspace{0.25em}
\resizebox{\linewidth}{!}{%
\begin{tabular}{@{}lccccccc@{}}
\toprule
Metric / Weakness & \textbf{R0} (v3.7.1) & \textbf{R1} & \textbf{R2} & \textbf{R3} & \textbf{R4} & \textbf{R5 (Rainbow~2)} & \textbf{R6} \\
\midrule
Published same-seed \& chosen-prefix witnesses collide? & \textcolor{Red}{\textbf{YES}} & \textbf{NO} & \textbf{NO} & \textbf{NO} & \textbf{NO} & \textcolor{Teal}{\textbf{NO}} & \textbf{NO} \\
\code{W1}: Theorem~\ref{thm:pairsum} 1-block state merges (of $1000$, $w=64$) & \textcolor{Red}{$1000$} & $0$ & \textcolor{Red}{$1000$} & $0$ & $0$ & \textcolor{Teal}{$\mathbf{0}$} & $0$ \\
\code{W1}: Distinct $(\sigma_1,\sigma_2)$ reached in 1 block ($w=4$, max $256$) & \textcolor{Red}{$1$} & $105$ & \textcolor{Red}{$1$} & $105$ & $105$ & \textcolor{Teal}{$\mathbf{105}$} & $105$ \\
\code{W1b}: $w=12$ DDT max ray bias $M_r$ / quotient $S_r$ & \textcolor{Red}{$11.0\times / 145\times$} & $3.95\times / 85.5\times$ & $8.98\times / 126\times$ & $3.67\times / 96.1\times$ & $3.67\times / 96.1\times$ & \textcolor{Teal}{$\mathbf{3.67\times / 96.1\times}$} & $3.70\times / 95\times$ \\
\code{W2}: Cross-pair avalanche after 1 $\mixA$ ($P_0 \leftrightarrow P_1$) & \textcolor{Red}{$0.000$} & \textcolor{Red}{$0.000$} & $0.229$ & $0.229$ & $0.229$ & \textcolor{Teal}{$\mathbf{0.229}$} & $0.215$ \\
\code{W2}: Cross-pair avalanche after 2 $\mixA$s ($P_0 \leftrightarrow P_1$) & \textcolor{Red}{$0.000$} & \textcolor{Red}{$0.000$} & $0.500$ & $0.500$ & $0.500$ & \textcolor{Teal}{$\mathbf{0.500}$} & $0.372$ \\
\code{W3}: $w=4$ finalizer reachable image (of $65{,}536$) & \textcolor{Red}{$4{,}007$ ($6.1\%$)} & \textcolor{Red}{$4{,}007$ ($6.1\%$)} & $41{,}766$ ($63.7\%$) & $41{,}766$ ($63.7\%$) & $41{,}390$ ($63.2\%$) & \textcolor{Teal}{$\mathbf{41{,}390}$ ($\mathbf{63.2\%}$)} & $41{,}395$ ($63.2\%$) \\
\code{W3}: $w=64$ dead-word $h_0''$ \& $(h_0,h_1)$ live (of $2000$) & \textcolor{Red}{$0 / 2000$} & \textcolor{Red}{$0 / 2000$} & $2000 / 2000$ & $2000 / 2000$ & $2000 / 2000$ & \textcolor{Teal}{$\mathbf{2000 / 2000}$} & $2000 / 2000$ \\
\code{W4}: Cross-width output prefixing (of $200$) & \textcolor{Red}{$200 / 200$} & \textcolor{Red}{$200 / 200$} & \textcolor{Red}{$200 / 200$} & \textcolor{Red}{$200 / 200$} & \textcolor{Red}{$200 / 200$} & \textcolor{Teal}{$\mathbf{0 / 200}$} & \textcolor{Red}{$200 / 200$} \\
\code{W5}: $\Delta\text{seed}=2^{63}$ IV Hamming weight (of 256 bits) & \textcolor{Red}{$4.0$} & \textcolor{Red}{$4.0$} & \textcolor{Red}{$4.0$} & \textcolor{Red}{$4.0$} & \textcolor{Red}{$4.0$} & \textcolor{Teal}{$\mathbf{76.7}$} & \textcolor{Red}{$4.0$} \\
\code{W5}: $w=8$ related-seed collision transfer rate & \textcolor{Red}{$9.5\%$} & \textcolor{Red}{$9.5\%$} & \textcolor{Red}{$9.5\%$} & \textcolor{Red}{$9.5\%$} & \textcolor{Red}{$9.5\%$} & \textcolor{Teal}{$\mathbf{0.5\%}$} & \textcolor{Red}{$9.5\%$} \\
\bottomrule
\end{tabular}%
}
\end{table}

\begin{table}[ht!]
\centering
\caption{\code{SMHasher3} verification codes, small-key latency ($[1,31]$\,B average cycles/hash), bulk throughput ($256$\,KiB aligned bytes/cycle and GiB/s at $3.5$\,GHz), and statistical suite results on 32-vCPU AMD EPYC Milan (\code{c2d-highcpu-32}).}
\label{tab:smhasher3-ladder}
\vspace{0.25em}
\resizebox{\linewidth}{!}{%
\begin{tabular}{@{}llcccc@{}}
\toprule
Rung & Widths (LE Verification Code) & 64b / 128b / 256b Small Keys & Bulk Throughput ($256$\,KiB) & Sanity (LE/BE) & Statistical Suite (\code{All}) \\
\midrule
\textbf{R0} (v3.7.1) & \code{ED7533D3} / \code{FF03173F} / \code{65F4A210} & $50.53$ / $56.94$ / $87.13$\,cyc & $1.47$\,B/c ($4.79$\,GiB/s) & 6/6 PASS & PASS (64/128/256) \\
\textbf{R1} (Asym.\ Load) & \code{095F02F9} / \code{FBDC549E} / \code{570AB08D} & $49.87$ / $53.68$ / $85.85$\,cyc & $1.47$\,B/c ($4.79$\,GiB/s) & 6/6 PASS & PASS (64/128/256) \\
\textbf{R2} (Feistel $\mixA$) & \code{E02D3EC1} / \code{5CF6A260} / \code{A5D67BE2} & $56.99$ / $67.22$ / $112.11$\,cyc & $1.39$\,B/c ($4.52$\,GiB/s) & 6/6 PASS & PASS (64/128/256) \\
\textbf{R3} (\code{R1+R2}) & \code{E2B79318} / \code{E09013C2} / \code{D32AAA3D} & $56.35$ / $64.65$ / $112.19$\,cyc & $1.44$\,B/c ($4.69$\,GiB/s) & 6/6 PASS & PASS (64/128/256) \\
\textbf{R4} (\code{R3}+4-Lane) & \code{CA69A2EB} / \code{8CD2FC6F} / \code{21A65AE5} & $56.45$ / $82.62$ / $130.20$\,cyc & $1.39$--$1.44$\,B/c ($4.51$--$4.70$\,GiB/s) & 6/6 PASS & PASS (64/128/256) \\
\textbf{R5 (Rainbow~2)} & \code{B6EB54B1} / \code{BA2E84CC} / \code{4F7D97BE} & $\mathbf{59.01}$ / $\mathbf{81.93}$ / $\mathbf{130.98}$\,\textbf{cyc} & $\mathbf{1.39}$--$\mathbf{1.44}$\,\textbf{B/c} ($\mathbf{4.51}$--$\mathbf{4.70}$\,\textbf{GiB/s}) & \textbf{6/6 PASS} & \textbf{PASS (64/128/256)} \\
\textbf{R6} (Cyclic $\mixA$) & \code{F502226B} / \code{5002C45E} / \code{24832643} & $45.66$ / $55.05$ / $95.61$\,cyc & $1.41$--$1.42$\,B/c ($4.61$--$4.62$\,GiB/s) & 6/6 PASS & PASS (64/128/256) \\
\bottomrule
\end{tabular}%
}
\end{table}

\subsection{Design Rationale: Why \code{R5} Is Selected as Rainbow~2}
\label{sec:rationale}

Comparing the seven rungs across Tables~\ref{tab:algebraic-ladder} and~\ref{tab:smhasher3-ladder} clarifies the exact engineering trade-offs:
\begin{enumerate}
  \item \textbf{Why not \code{R1} alone?} \code{R1} is remarkable for breaking the Theorem~\ref{thm:pairsum} pair-sum coset collapse at zero speed cost ($49.87$\,cycles/hash small keys, $4.79$\,GiB/s bulk). However, \code{R1} leaves $\mixA$ pair-separable (\code{W2}) and retains the 192-bit finalization bottleneck (\code{W3}), where word $h_0''$ is $100\%$ dead in 64-bit mode and the $w=4$ finalizer reaches only $6.1\%$ of states.
  \item \textbf{Why not \code{R3} or \code{R4}?} \code{R3} and \code{R4} fix the streaming and finalization core (\code{W1}--\code{W3}) at $4.61$--$4.69$\,GiB/s, but retain the linear additive IV (\code{W4}--\code{W5}), leaving \code{rainbow-64}$(m)$ as a trivial prefix of \code{rainbow-128/256}$(m)$ and leaving the top-bit seed differential $\Delta\text{seed}=2^{63}$ with a 4-bit IV Hamming weight.
  \item \textbf{Why not \code{R6}?} \code{R6} achieves the lowest small-key latency on the ladder ($45.66$\,cycles/hash) by placing $b \mathrel{\oplus}= c,\ d \mathrel{\oplus}= a$ prior to the second multiply in $\mixA$ and cyclically renaming the output registers. However, its 2-round cross-pair avalanche ($0.372$) is slightly below the full $0.500$ diffusion of the post-multiply Feistel cross-coupling in \code{R2}--\code{R5}, and without \code{R5}'s IV it leaves \code{W4} and \code{W5} open.
  \item \textbf{Why \code{R5} wins as Rainbow~2 (\code{rainbow2}):} \code{R5} is the \textbf{only rung on the ladder that eliminates every structural weakness identified in our cryptanalytic campaigns (\code{W1}, \code{W1b}, \code{W2}, \code{W3}, \code{W4}, and \code{W5}) simultaneously}. Its bulk throughput ($1.41$--$1.44$\,bytes/cycle, $4.60$--$4.69$\,GiB/s) remains within $2\%$--$4\%$ of unmodified Rainbow~1 ($1.47$\,bytes/cycle, $4.79$\,GiB/s), because the domain-separated rotated-multiply IV and the rotated four-lane output squeeze execute only once per message outside the inner 16-byte block loop. We register \textbf{Rainbow~2} (\code{rainbow2}) as a separate hash family so existing Rainbow~1 test vectors remain untouched while downstream users and \code{SMHasher3} gain a hardened 64/128/256-bit family whose wider digests deliver full-state mixing across all 256 state bits.
\end{enumerate}

\section{Reproducibility}
\label{sec:repro}

All sources, witnesses, hardening ladder variants, and machine-readable results are published in \code{research/rainbow/}, \code{research/rainbow-spectrum/}, \code{research/rainbow2/}, and \code{src/rainbow2.cpp}. From the repository root:

\begin{lstlisting}[style=code]
# 1. Rainbow v3.7.1 same-seed and chosen-prefix collisions & verification
clang++ -std=c++20 -O3 -march=native -o research/rainbow/bin/collision_pair_sum \
    research/rainbow/collision_pair_sum.cpp
clang++ -std=c++20 -O3 -march=native -o research/rainbow/bin/chosen_prefix_collision \
    research/rainbow/chosen_prefix_collision.cpp
clang++ -std=c++20 -O3 -march=native -o research/rainbow/bin/replay_production \
    research/rainbow/replay_production.cpp
python3 research/rainbow/verify_pair_sum.py research/rainbow/collision-result.json
python3 research/rainbow/verify_pair_sum.py research/rainbow/chosen-prefix-result.json

# 2. Rainbow 2 hardening ladder (R0..R6) generation, ASan/UBSan & algebraic verification
python3 research/rainbow2/build_variants.py
python3 research/rainbow2/algebraic_ladder_eval.py
python3 research/rainbow2/ddt_and_seed_eval.py
\end{lstlisting}

\section{Conclusion}
\label{sec:conclusion}

Rainbow v3.7.1's message schedule preserves two 64-bit pair sums that an attacker can move at will within their coset, its primary mixer $\mixA$ never crosses the two 128-bit halves of the state, and its two-lane output squeeze leaves 64 bits of state dead in 64-bit finalization. Together these properties reduce a nominally 256-bit collision problem to two independent 64-bit birthday searches, yielding practical same-seed ($2^{32.28}$) and chosen-prefix ($2^{34.80}$) collisions against the full unreduced hash across all three digest widths simultaneously.

Using these exact structural findings as design specifications, we constructed a seven-rung hardening ladder (\code{R0}--\code{R6}) and evaluated it through algebraic verifiers and 32-core \code{SMHasher3} cloud campaigns. The resulting algorithm, \textbf{Rainbow~2} (\code{R5}), combines asymmetric rotated block injection, Feistel cross-coupled $\mixA$, a rotated four-lane output squeeze, and a domain-separated rotated-multiply IV to eliminate all five classes of structural weaknesses (\code{W1}--\code{W5}) while sustaining $4.60$--$4.69$\,GiB/s bulk throughput and passing the full \code{SMHasher3} suite across 64-, 128-, and 256-bit digests.

\begin{thebibliography}{99}

\bibitem{rain-source}
DOSAYGO Corporation.
\newblock \emph{rain: Rainbow and Rainstorm hash functions}, \code{src/rainbow.cpp}, version 3.7.1.
\newblock \url{https://github.com/DOSAYGO-STUDIO/rain}

\bibitem{rain-readme}
DOSAYGO Corporation.
\newblock \emph{rain README}: ``Rainbow \dots General-purpose, non-crypto hash \dots Not designed as cryptographic.''

\bibitem{rainstorm-cryptanalysis}
Cris.
\newblock \emph{A $2^{31.9}$ Collision Attack on Reduced Rainstorm-256: Programmed Collisions, Rebound Corridors, and the Full-Schedule Boundary}.

\bibitem{vow1999}
P. C. van Oorschot and M. J. Wiener.
\newblock \emph{Parallel Collision Search with Cryptanalytic Applications}.
\newblock Journal of Cryptology, 12(1):1--28, 1999.

\bibitem{joux2004}
A. Joux.
\newblock \emph{Multicollisions in Iterated Hash Functions. Application to Cascaded Constructions}.
\newblock CRYPTO 2004, LNCS 3152, pp. 306--316.

\bibitem{kelseyschneier2005}
J. Kelsey and B. Schneier.
\newblock \emph{Second Preimages on $n$-Bit Hash Functions for Much Less than $2^n$ Work}.
\newblock EUROCRYPT 2005, LNCS 3494, pp. 474--490.

\bibitem{kelseykohno2006}
J. Kelsey and T. Kohno.
\newblock \emph{Herding Hash Functions and the Nostradamus Attack}.
\newblock EUROCRYPT 2006, LNCS 4004, pp. 183--200.

\bibitem{stevens2007}
M. Stevens, A. K. Lenstra, and B. de Weger.
\newblock \emph{Chosen-Prefix Collisions for MD5 and Colliding X.509 Certificates for Different Identities}.
\newblock EUROCRYPT 2007, LNCS 4515, pp. 1--22.

\bibitem{merkle1989}
R. C. Merkle.
\newblock \emph{One Way Hash Functions and DES}.
\newblock CRYPTO 1989, LNCS 435, pp. 428--446.

\bibitem{smhasher3}
Various contributors.
\newblock \emph{SMHasher3: Hash function quality and speed test suite}.

\end{thebibliography}

\end{document}
